\documentclass[../../main.tex]{subfiles}
\begin{document}
\section{Route C: the deterministic moment-map/Haar program, audited frontier}
\label{sec:moment-map-cmh}

\subsection*{Route gateway}
\addcontentsline{toc}{subsection}{Route gateway}

\begin{description}[leftmargin=0pt,style=nextline,itemsep=0.3em]
  \item[1. Thesis.] Bound one deterministic quantity, the canonical moment-Hessian constant
  $\CMH$, and pass it to an affine Poincar\'e inequality. No stochastic localization anywhere.

  \item[2. Bridge to KLS.] The strongest bridge of any route here, and the one place this
  document proves an implication \emph{into} the neighbourhood of KLS rather than assuming one:
  \[
    \mathrm{CMH}(4)
    \ \xRightarrow[\ \text{proved: Thm.~\ref{thm:cmh-implies-affine-poincare}}\ ]{}\
    \CPaff\le\CMH
    \ \xrightarrow[\ \texttt{ass:uniform-cmh-approximants}\ ]{\ \text{open premise}\ }\
    \text{all log-concave }\mu .
  \]
  Two cautions belong here rather than in a footnote, because they change how the diagram
  should be read. $\mathrm{CMH}(4)$ is \emph{not} a reformulation of
  Conjecture~\ref{conj:kls}: by Proposition~\ref{prop:cmh-hodge} it additionally demands
  control of a solenoidal excess that vanishes identically in dimension one, so it may be
  strictly stronger --- equivalence at constant~$4$ is open
  (Corollary~\ref{rem:cmh-stronger-than-kls}). And the passage to arbitrary log-concave limits
  is certified only \emph{conditionally}: Question~\ref{q:cmh-approximation} is proved under
  the separate open premise \texttt{ass:uniform-cmh-approximants}.

  \item[3. Established inputs.] From the literature: the fixed-matrix Hessian bound
  (Theorem~\ref{thm:letwin-moment-map}) and the moment-Hessian trace bound
  (Theorem~\ref{thm:chen-klartag-moment-hessian}); moment-potential regularity from
  \cite{BermanBerndtsson2013RealMA} and the Stein identity from
  \cite{Fathi2019SteinMomentMaps}. Certified here: the endpoint's operator data
  (Definition~\ref{def:cmh}) and everything in
  Sections~\ref{sec:cmh-normalization}--\ref{sec:cmh-exact-cases} that carries a
  \emph{certified} standing.

  \item[4. Main advance so far.] The endpoint was made precise and then computed. It is given
  operator data rather than left as a slogan (Definition~\ref{def:cmh}); it is proved to
  dominate the affine Poincar\'e constant (Theorem~\ref{thm:cmh-implies-affine-poincare}); its
  Hodge content is identified, separating an affine channel from a solenoidal excess
  (Proposition~\ref{prop:cmh-hodge}); and it is evaluated \emph{exactly} on the line, on
  products, and on every log-concave Dirichlet law (Section~\ref{sec:cmh-exact-cases}). That
  last is a genuine theorem about a nontrivial family, not a consistency check.

  \item[5. Exact bottleneck.] The route splits, and the split is the point. A
  \emph{construction} layer: \texttt{q:mm-invariant-lift}, deriving the target-flat
  Schur--Piola multiplier lift invariantly, and \texttt{q:mm-square-root-commutator},
  controlling the full Haar sum of $[N^{1/2},K_M]$ errors without double-spending the positive
  reservoir. A \emph{falsification} layer: gate zero (Conjecture~\ref{conj:gate-zero}) and the
  solenoidal perturbation test \texttt{q:cmh-solenoidal-perturbation}. The route can be closed
  by the first pair or killed by the second.

  \item[6. Failed variants and obstructions.] Proposition~\ref{prop:letwin-not-gate-zero} is
  the decisive negative result: no matrix-moment argument supplies gate zero, so the
  fixed-matrix estimate cannot be leveraged into the linear sector. Gate zero is itself an
  instance of the average-versus-uniform pattern of Section~\ref{sec:kls-remaining} --- it asks
  $\lmax(\E H^2)\le4$ where only $\Tr(\E H^2)\le2n$ is known --- so this route relocates that
  difficulty rather than escaping it, and Proposition~\ref{prop:letwin-not-gate-zero} shows the
  relocation is not free. A separate saturation risk is recorded in
  Section~\ref{sec:cmh-exact-cases}.

  \item[7. Completion or reopening criterion.] The route closes if the construction layer is
  discharged \emph{and} \texttt{ass:uniform-cmh-approximants} is proved, or replaced by the
  weaker recovery envelope \texttt{ap:cmh-recovery-envelope}. It is refuted if gate zero is
  false, or if the exact second variation through the Hodge split at the saturating
  one-sided-exponential product is negative. Both falsifiers are cheaper than the
  construction, which is why they are stated as gates.

  \item[8. Technical reading guide.] Read \emph{out of order}: the normalization layer,
  Section~\ref{sec:cmh-normalization}, is logically prior and should be read first, then the
  exact cases in Section~\ref{sec:cmh-exact-cases}, then this construction layer. The
  moment-map family survey is Section~\ref{sec:family-moment-map}. None of
  Appendices~\ref{sec:notation}--\ref{sec:models} is used by this route.
\end{description}

This section records a second, deterministic use of the moment map. Its aim is to extend
Theorem~\ref{thm:letwin-moment-map} from constant matrices to multiplier fields selected by an
arbitrary test function. The route is live but unproved. Displayed formal identities below are
not proof-certified; current entry points are the approaches of
\texttt{research/program/portfolio.yaml}, and each approach's \texttt{objective} states what
closing it would deliver.

Route C is organized in two layers. This section is the \emph{construction} layer: Haar
compression, Schur--Piola transport, and the square-root commutator that the programme must
control. Sections~\ref{sec:cmh-normalization} and \ref{sec:cmh-exact-cases} are the
\emph{normalization} layer, added later and logically prior: they fix the endpoint's operator
data, prove that it implies the affine Poincar\'e inequality, compute it exactly on the line, on
products, and on every log-concave Dirichlet law, and isolate the two falsification tests the
route now admits. The two layers share the target $\mathrm{CMH}(4)$ and almost no machinery; a
reader starting Route C should read \S\ref{sec:cmh-normalization} first.

To avoid a collision with the stochastic quantities $H_t,S_t,K_t$ and $A_t$ used earlier, all
objects in this section are stationary moment-map objects: $H=D^2\phi$ is the Hessian metric,
$N$ is a weighted elliptic operator, and $K_M$ is a compressed multiplier.

\begin{program}[Deterministic CMH route]
\label{prog:cmh-route}
First make the CMH endpoint and its KLS reduction precise. Then derive the invariant multiplier
lift, control the resulting square-root commutators with the full positive reservoir, and sum the
complete Haar tree without nodewise positivity or duplicated slack.
\end{program}

\subsection{Endpoint and first audit gate}

The originating exploration proposes a covariance--moment--Hessian estimate of the schematic
form
\begin{equation}
  \norm{\Sigma^{-1/2}H\nabla g}_2^2
  \le4\norm{-Lg}_2^2.
  \label{eq:cmh4-schema}
\end{equation}
It reports that a divergence-duality argument would then give $\CP(\mu)\le4$, the sharp plausible
constant because a standard centered one-sided exponential has Poincar\'e constant $4$.

\begin{remark}[CMH normalization and endpoint reduction]
\label{q:cmh-normalization}
Define $\Sigma$, $L$, the underlying $L^2$ space, the admissible class of $g$, and every inverse
in \eqref{eq:cmh4-schema} for the regular moment-map class. Then prove that the resulting
estimate $\CMH(\mu)\le C$ implies $\CPaff(\mu)\le C$ on that class.
\end{remark}

This regular-class question is logically prior to the commutator calculation, and it is now
answered.
Definition~\ref{def:cmh} fixes $\Sigma$ as the covariance, $L$ as the Stein generator
$\Div_\mu(H\nabla\,\cdot\,)$, the $L^2$ space as $L^2(\mu)$, the admissible class as
$\Dom(\Aop)$, and every inverse as the pseudoinverse on $(\ker\Aop)^\perp$;
Theorem~\ref{thm:cmh-implies-affine-poincare} then proves
$\CPaff(\mu)\le\CMH(\mu)$ by a single Cauchy--Schwarz step, with a spectral truncation in place
of an assumed gap.

\begin{assumption}[Uniform CMH control on regular approximants]
\label{ass:uniform-cmh-approximants}\klsstatus{ass:uniform-cmh-approximants}
There is a universal $C<\infty$ such that, for every centered log-concave law $\mu$, the
centered Gaussian-convolution, Gaussian-tilt, and growing-ball regular moment-map approximants
$\mu_k$ satisfy $\sup_k\CMH(\mu_k)\le C$.
\end{assumption}

\begin{question}[Approximation closure for CMH]
\label{q:cmh-approximation}\klsstatus{q:cmh-approximation}
Under Assumption~\ref{ass:uniform-cmh-approximants}, establish a closure argument under which the
affine Poincar\'e inequalities pass to $\mu$, uniformly under isotropic normalization and
affine-support degeneration.
\end{question}

The approximation closure is now certified conditionally. Centered Gaussian-convolution,
Gaussian-tilt, and growing-ball approximants converge in ambient $W_2$, and the affine
Poincar\'e inequality passes with no loss on the intrinsic closed covariance-form domain,
including proper affine-support degeneration. The uniform premise
$\sup_k\CMH(\mu_k)\le C$ remains unresolved; no bound on $\CMH$ is supplied by the closure
argument. Two consequences of the now-explicit normalization bear directly on the rest of this
section.

First, the reduction is not known to be an equivalence. Proposition~\ref{prop:cmh-hodge} splits
the CMH numerator into an affine Poincar\'e part and a nonnegative solenoidal excess, which
vanishes identically in dimension one under the no-flux convention.  Thus CMH contains an
additional channel that KLS does not directly control; no separating log-concave measure is
currently known, so strict non-implication is not asserted
(Corollary~\ref{rem:cmh-stronger-than-kls}). The programme below should be read as pursuing a
sufficient condition of unknown truth value, not as a proved reformulation of the conjecture.

Second, the endpoint now has a cheap necessary condition. Testing $\mathrm{CMH}(4)$ on linear
functions gives gate zero, $\E[H\Sigma^{-1}H]\preceq4\Sigma$
(Conjecture~\ref{conj:gate-zero}), and Proposition~\ref{prop:letwin-not-gate-zero} shows by an
exact countermodel that Theorem~\ref{thm:letwin-moment-map} does not imply it through matrix
algebra alone. The obstruction there is the \emph{static} commutator
\eqref{eq:static-commutator}, $\Tr(B^2H^2)=\Tr(BHBH)+\tfrac12\norm{[B,H]}_{\HS}^2$ --- the
finite-dimensional shadow of the square-root commutator of
Question~\ref{q:mm-square-root-commutator}. A programme that controls
$[N^{1/2},K_M]$ must in particular control $\E\norm{[B,H]}_{\HS}^2$; conversely, refuting gate
zero would close this route without any Haar-tree analysis at all.

\subsection{Haar compression and the commutator error}

On the mean-zero subspace, formally set
\[
  N=D_\nu^*H^{-1}D,
  \qquad
  Q_M=D_\nu^*H^{-1/2}MH^{-1/2}D,
\]
and
\[
  R=H^{-1/2}DN^{-1/2},
  \qquad K_M=R^*MR.
\]
On a common core, $R^*R=I$ and
\[
  Q_M=N^{1/2}K_MN^{1/2}.
\]
For normalized Haar contrast multipliers $M_S$, the formal error is therefore
\begin{equation}
  e_S(u)=Q_{M_S}u-K_{M_S}Nu
  =[N^{1/2},K_{M_S}]N^{1/2}u.
  \label{eq:mm-commutator-error}
\end{equation}
The reported complete-tree Bessel deficit is
\begin{equation}
  \mathfrak B(h)
  =\left(1-\frac1n\right)\norm{h}^2-\sum_S\norm{K_{M_S}h}^2
  =\sum_S\norm{(I-RR^*)M_SRh}^2.
  \label{eq:mm-bessel-deficit}
\end{equation}
The originating exploration reports rotated examples with negative individual sibling deficits
and Laguerre near-extremizers that consume nearly all descendant slack. Pending persisted
witnesses, these are route-level retractions rather than ledger refutations; they rule out
nodewise positivity and fixed fractional slack allocation within the originating analysis.

\subsection{Schur--Piola transport and local algebra}

For a $1+d$ split write
\[
  H=\begin{pmatrix}h&b^\top\\ b&C\end{pmatrix},\qquad
  v=C^{-1}b,\qquad s=h-b^\top C^{-1}b,\qquad \Delta=\det C,
\]
and $X=\partial_r-v\cdot\nabla_q$. Block algebra, Piola's identity, and Hessian compatibility
formally give
\[
  \operatorname{cof}H
  =\Delta\begin{pmatrix}1&-v^\top\\-v&vv^\top+sC^{-1}\end{pmatrix},
  \qquad
  \operatorname{div}(\Delta X)=0,
\]
\[
  Xv=C^{-1}\nabla_qs,
  \qquad XC=J^\top C=CJ,
  \qquad \Tr(C^{-1}XC)=\operatorname{div}_qv.
\]
If $z=\nabla_q\phi$ is the child target coordinate, then $Xz=0$; for
$x=\nabla\phi$, $Xx_\perp=0$ and $Xx_0=s$. Hence $s^{-1}X$ is the pullback of
$\partial_{x_0}$. In particular, the trace/conformal derivative is constrained by the Schur
connection and is not a free scalar mode.

At a Schur-normal $1+2$ point, let $T\in\operatorname{Sym}_2$, $a,g\in\R^2$, let
$R_\perp e_1=e_2$ and $R_\perp e_2=-e_1$, and use
$\operatorname{sym}(u\otimes v)=\tfrac12(u\otimes v+v\otimes u)$. Put
\[
  B_i=\operatorname{sym}(g\otimes R_\perp^\top e_i).
\]
The fixed-target coordinate form expands exactly as
\begin{align*}
Q(a,T,g)
&=8|a|^2+\frac12\sum_i\norm{\{B_i,T\}}_{\HS}^2-4a\cdot R_\perp Tg\\
&=8\left|a-\frac14R_\perp Tg\right|^2
  +\frac32|Tg|^2+\frac14(\Tr T)^2|g|^2\ge0.
\end{align*}
An independent coordinate expansion checked this identity. What remains open is the invariant
identification of the tensorial lift actually generated by the global operator, and a reduction
covering all higher-dimensional splits. Positivity in $1+2$, $1+3$, and $2+2$ alone would not be
a dimension-free theorem.

\begin{question}[Invariant lift and all-split reduction]
\label{q:mm-invariant-lift}\klsstatus{q:mm-invariant-lift}
Derive the multiplier transport in a target-flat Schur--Piola frame without choosing a moving-frame
coefficient by hand. Decompose the resulting Codazzi shape terms into controlled irreducible
pieces for arbitrary split dimensions, and match each negative component with a retained
Letwin/Monge--Amp\`ere square.
\end{question}

\subsection{Two solenoidal channels}

Let $S(z)=\E[C\mid z]$ be the inherited child block of the parent Stein kernel and let $K(z)$ be
the canonical child moment-map Stein kernel. Both solve the same Stein equation, so
\[
  \mathsf D=S-K,
  \qquad \partial_\beta(\rho\mathsf D_{\alpha\beta})=0.
\]
In a two-dimensional child, subject to topology and boundary conditions,
\[
  \rho\mathsf D=R_\perp^\top D^2\lambda R_\perp
\]
is the Airy representation. For a parent datum $h$, the reported canonical flux residual has the
form
\begin{equation}
  r_h=r_{\mathrm{cond}}+(I-\Pi_K)\mathsf D\nabla v.
  \label{eq:mm-hodge-residual}
\end{equation}
Thus a conditional-flux channel and a canonical-versus-inherited Stein-kernel channel must be
controlled separately. The earlier guessed local conserved vector collapsed these two projections
and was withdrawn.

\subsection{The best identified global bottleneck}

The cubic symbol does not cancel in general. With $A_M=H^{-1/2}MH^{-1/2}$, the originating
calculation gives
\[
  C_3(\xi)=2\left[(H^{-1}\xi)^r(\partial_rA_M)(\xi,\xi)
  -(A_M\xi)^r(\partial_rH^{-1})(\xi,\xi)\right].
\]
It is proposed to equal
\[
  2(H^{-1}\xi)^k\xi^\top H^{-1/2}[\Omega_k,M]H^{-1/2}\xi,
  \qquad
  \Omega_k=\frac12\left(H^{-1/2}\partial_kH^{1/2}
  -(\partial_kH^{1/2})H^{-1/2}\right),
\]
exhibiting eigenframe rotation rather than conformal variation. The square-root bridge is the
standard resolvent formula
\[
  [N^{1/2},K_M]
  =\frac1\pi\int_0^\infty t^{1/2}(N+t)^{-1}
  [N,K_M](N+t)^{-1}\dd t,
\]
provided all forms and domains are justified.

\begin{question}[Global square-root commutator]
\label{q:mm-square-root-commutator}\klsstatus{q:mm-square-root-commutator}
Define a retained nonnegative remainder $\mathfrak R_{\mathrm{Letwin}}(u)$ containing every
available variable-multiplier, Codazzi, Monge--Amp\`ere, corrector, and descendant square, and
prove over the complete Haar tree that
\[
  2\sum_S\operatorname{Re}\langle K_{M_S}Nu,e_S(u)\rangle
  +\sum_S\norm{e_S(u)}^2
  \le \mathfrak B(Nu)+\mathfrak R_{\mathrm{Letwin}}(u)
\]
with a dimension-free constant and stable operator domains. Then show that the complete-tree
estimate supplies the global analytic input required by the fully defined CMH target of
Question~\ref{q:cmh-normalization}.
\end{question}

This is the best identified bottleneck, not yet the sole formal gap. The invariant lift, all-split
reduction, Hodge boundary conditions, and one-edge flux identity must be closed alongside it; the
endpoint duality is no longer among them, having been discharged by
Theorem~\ref{thm:cmh-implies-affine-poincare}. Note also that
Proposition~\ref{prop:letwin-not-gate-zero} places a floor under the difficulty: even the
constant-multiplier, static specialization of the estimate below is not available from
\eqref{eq:letwin-matrix} by algebra, so no argument here may treat the commutator as a
lower-order correction. The originating regression analysis indicates that no scalar, conformal-only, or
nodewise shortcut survives its current models. Persisting those witnesses is still required, but
the resulting proof design is necessarily global and matrix-coupled.
\end{document}
